\section{Code quality, tests, README and packaging}
\label{sec:quality}

Five points are for the source code: it runs with no errors or warnings (2), it is commented (1)
and it is well structured (2).

\subsection{Code rules}
\label{sec:coderules}

\textbf{Both tiers.}
\begin{enumerate}
  \item Every file opens with a comment or docstring that says what it does and which requirement,
    scope item or endpoint it serves (for example ``E8, M-03'').
  \item Comments explain why, not what: why validation errors answer 400, why hypothetical rows
    share one time, why the Assistant's text is never streamed. Comments that restate the code are
    left out.
  \item No number or message key is typed into code. Numbers come from \fp{settings.json}; the
    front end gets them from E2.
  \item No dead code, no commented-out experiments, no leftover console or print output.
  \item Names say what a thing is in the product's words: address, transaction, flag, method,
    what-if.
\end{enumerate}

\textbf{Back end}, as in Assignment 2: type hints throughout; a docstring with arguments, units,
return value and errors on every public function; \code{pathlib} for every path; routers thin,
services free of HTTP, the copied package untouched (Section~\ref{sec:architecture}).

\textbf{Front end.} TypeScript in strict mode and no \code{any}; one component per file with
typed props; views compose components and never call Axios or Plotly directly; styles use the
design tokens; ESLint and Prettier run before every commit.

\subsection{Tests}
\label{sec:tests}

{\small
\begin{longtable}{@{}L{2.5cm}L{2.2cm}L{9.4cm}@{}}
\toprule
\textbf{Level} & \textbf{Tool} & \textbf{What it covers} \\
\midrule
\endhead
Back end, unit & pytest & The address and transaction-ID checkers against the shared vectors; the what-if rules; the batched what-if against single calls; the CSV formula guard; the watchlist store; settings loading \\
Back end, API & pytest with FastAPI's test client & Every endpoint: its success, and each status in its row of Section~\ref{sec:endpoints}. A wrong type posted to every POST and PUT answers 400, never 422 or 500 \\
Parity with Assignment 2 & pytest & For all 28 addresses, E5, E6 and E7 together rebuild exactly what \code{Predictor.score_address} returns; each method's flags equal \fp{scores.csv}; E8 equals \code{score_what_if} \\
Wording & pytest; Vitest & Every string in every response for all 28 addresses and a grid of what-ifs passes \code{check_outputs}; every rendered view and state has no banned word; every flagged item has the limitation line \\
Integrity & pytest & Every copied file matches its hash in \fp{A2_VENDOR.json}; the installed library versions equal the pins \\
API documentation & pytest & The OpenAPI schema builds; every route has a response model, a summary and an example; the committed copy of \fp{openapi.json} is current \\
Front end, unit & Vitest & The address checker and the what-if rules against the shared vectors; formatters; the router; the hooks (debounce, cancellation, polling) \\
Front end, component & Vitest and Testing Library & Every row of the state catalogue; every chart's states and table view; the what-if form's feedback; dialogs hold and return focus \\
Front end, journeys & Vitest & J1 to J6 on the in-browser mock, whose answers are fixtures recorded from the real API \\
Accessibility & axe in Vitest; by hand & No automatic violation on any view; the keyboard run and the checks of Section~\ref{sec:a11y} \\
Browser (S-11) & Playwright & J1 to J6 against both real servers at 390, 820 and 1280 pixels, in Chromium, Firefox and WebKit; no sideways scrolling; an empty browser console; the report's screenshots \\
Performance (S-12) & \code{measure.py}; a browser script & The budgets of Section~\ref{sec:budgets} \\
By hand & \fp{docs/TEST_LOG.md} & The journeys on real devices and browsers: date, build, tester, device, result, issue \\
Clean install & A teammate & The README, from the zip, on Windows 11 and on macOS (AT-11) \\
\bottomrule
\end{longtable}
}

\textbf{Rules.} \code{python -m pytest -q} and \code{npm test} both pass before every commit, and
the exit code is what counts. A bug gets a test that fails before the fix. The mock is used only by
tests; \code{npm run fixtures} refreshes its answers from the running back end, so it cannot drift
from the real contract. With S-10, the TypeScript types are generated from \fp{openapi.json} and a
check fails when they are stale, and GitHub Actions runs both suites, the type check and the build
on every push.

\subsection{README}
\label{sec:readme}

The brief asks for set-up instructions, the libraries needed, how to run both tiers, and how the
model is configured. The README follows this outline.

\begin{enumerate}
  \item \textbf{What this is}, in three sentences, with one screenshot and the folder map.
  \item \textbf{What you need}: Windows 10 or 11, macOS on Apple silicon, or Linux; conda
    (Miniforge or Miniconda); Node.js 22 or later; about 2 GB of disk; no GPU.
  \item \textbf{Start the back end}:
\begin{lstlisting}
cd backend
conda env create -f environment.yml
conda activate cos30049-a3
python -m pytest -q                      # expect: all passed
uvicorn app.main:app --port 8000         # expect: "Explorer ready: 28 addresses, 9 methods"
\end{lstlisting}
  \item \textbf{Start the front end}, in a second terminal:
\begin{lstlisting}
cd frontend
npm install
npm run dev                              # then open http://localhost:5173
\end{lstlisting}
  \item \textbf{Check it works}: open \code{http://localhost:8000/api/health}; open the
    application and pick the first address; the expected screen is shown.
  \item \textbf{A two-minute tour}: the journey of the video, in six steps, with the example
    address and a what-if to type.
  \item \textbf{The API}: the endpoint table and the link to \code{/docs}.
  \item \textbf{How the model is integrated and configured}: where the models and tables are, how
    they were copied from Assignment 2 (\fp{A2_VENDOR.json}), what \fp{settings.json} controls, the
    environment variables, and how to add the optional Assistant weights and check their SHA-256.
  \item \textbf{Tests}: the commands and what they cover.
  \item \textbf{Troubleshooting}: a port already in use; \code{conda} not found, or its terms of
    service not accepted; an older Node.js; the pip route (\code{pip install -r requirements.txt}
    on Python 3.11); a library version message at start; \code{set PYTHONUTF8=1} on Windows; the
    firewall prompt on first start; a blank page when the back end is not running.
  \item \textbf{Libraries and licences}, and the acknowledgement of AI use.
\end{enumerate}

\begin{team}
\textbf{The clean install (T44, milestone MS4).} A member who did not write the README follows it
from a trial zip on a machine that has never had the project: Windows 11 first, then macOS. They
note every place they had to guess and the time each step took. Every guess becomes a README fix.
Assignment 1's risk 9 is closed only when they finish without help and both terminals and the
browser console are free of errors and warnings.
\end{team}

\subsection{The zip}
\label{sec:packaging}

\code{python tools/package.py --tag v1.0} builds the zip from the git tag, never from the working
folder, and writes \fp{SHA256SUMS.txt} into it and beside it, as in Assignment 2.

\begin{center}
\small
\begin{tabularx}{\textwidth}{@{}L{5.0cm}Y L{2.0cm}@{}}
\toprule
\textbf{In the zip} & \textbf{Contents} & \textbf{Size}\est{} \\
\midrule
\fp{README.md}, \fp{LICENSES.md}, \fp{CHANGELOG.md} & Run guide; licences; changes & under 1 MB \\
\fp{backend/app/}, \fp{backend/tests/} & The FastAPI service and its tests & under 1 MB \\
\fp{backend/explorer/}, \fp{backend/lm/} & The model code used by the web application, copied from Assignment 2 & under 1 MB \\
\fp{backend/MODIFICATIONS.md}, \fp{A2_VENDOR.json} & What was changed for the web; where each file came from & -- \\
\fp{backend/models/} & The nine saved methods and the profile model & 31 MB \\
\fp{backend/data/processed/esplora/} & The snapshot's features, transactions, scores, TxGPT scores and manifest & 12 MB \\
\fp{backend/config/settings.json} & Every threshold and limit & -- \\
\fp{frontend/} & The React source, \fp{package.json} and \fp{package-lock.json} & 1 MB \\
\fp{shared/} & The validation test vectors & -- \\
\fp{docs/AI_USE_LOG.md} & The record behind the AI acknowledgement & -- \\
\midrule
\multicolumn{2}{@{}l}{Whole zip, compressed} & about 40 MB \\
\bottomrule
\end{tabularx}
\end{center}

\textbf{Left out:} \fp{node_modules/} (the brief says so), \fp{frontend/dist/}, \fp{.git/},
\fp{.env} files, \fp{backend/var/}, caches, the brief, this plan, the report's sources and the
working notes. The Assistant's weights (610 MB) are a download named in the README with their
SHA-256, if the tutor agrees (Section~\ref{sec:tutorq}).

\begin{done}
The zip unpacks into the tree of Section~\ref{sec:layout}; \fp{SHA256SUMS.txt} verifies every file;
on a clean machine the README's commands start both tiers and both test suites pass from the
unzipped folder with no git history.
\end{done}
